%----------------------------------------------------------------------------
\chapter{Felhasznált technológiák}
%----------------------------------------------------------------------------

A projekt megvalósításához számos modern technológiát alkalmaztam, amelyek lehetővé tették a hatékony fejlesztést és egy skálázható rendszer kialakítását.
A következőkben részletesen bemutatom az egyes technológiákat és azok szerepét a rendszerben.
\section{Szerver (Backend)}

\subsection{Python}

A Python \cite{Python} egy magas szintű, általános célú programozási nyelv, amelyet Guido van Rossum hozott létre 1991-ben. Könnyen olvasható szintaxisa és széles körű könyvtárkészlete miatt népszerű a webfejlesztés, adatfeldolgozás, gépi tanulás és sok más területen. A Python dinamikus típusosságot és objektumorientált programozást támogat, valamint lehetőséget biztosít aszinkron programozásra is. A projekthez a Python 3.13.3-as verzióját használtam.

\subsubsection{FastAPI}

A backend (szerver) fejlesztéséhez a FastAPI \cite{FastAPI} keretrendszert választottam, ami egy modern, gyors webkeretrendszer Python nyelven, amely lehetővé teszi REST-ful\footnote{REST: Representational State Transfer} API-k egyszerű fejlesztését. Fő előnye, hogy standard Python típusannotációkat használ és ezzel típusbiztonságot nyújt. Emellett aszinkron működést biztosít az endpointoknak az I/O-intenzív feladatoknál (pl. adatbázis lekérdesék, külső API hívások), gyors fejlesztést tesz lehetővé a beépített validáció és dokumentáció révén, és automatikus Swagger \cite{Swagger} dokumentációt is az endpointoknak\footnote{Endpoint: (magyarul végpont) Egy digitális cím (általában egy URL), amin keresztül a szerver-alkalmazás (API - Application Programming Interface) fogadja a külső kéréseket, és válaszként adatot szolgáltat.}.

\subsubsection{PostgreSQL}

Adatbázisnak a PostgreSQL-t \cite{PostgreSQL} választottam, ami egy nyílt forráskódú relációs adatbázis-kezelő rendszer. Azért esett erre a választás, mert a rendszer adatai között sok a kapcsolat, és ezeket relációs adatmodellben, idegen kulcsokkal könnyű volt leképzeni. A pénzmozgások miatt fontos volt az is, hogy az adatbázis tranzakciókezelése megbízható és visszavonható legyen egy tranzakció, ha valami hiba történt. A PostgreSQL ezt teljes mértékben támogatja (ACID\footnote{ACID: Atomicity, Consistency, Isolation, Durability - az adatbázis-tranzakciók négy alapvető tulajdonsága.} tulajdonságok).

Emellett a PostgreSQL-ben van néhány olyan adattípus, amit más adatbázisokban nem, vagy csak körülményesebben lehet használni, például a \texttt{JSONB} típus, amiben strukturálatlan adatot is el lehet tárolni, vagy a natív \texttt{UUID} és \texttt{ENUM} típusok.

\subsubsection{SQLAlchemy}

Az adatbázis-hozzáférés nem nyers SQL-lekérdezésekkel, hanem az SQLAlchemy \cite{SQLAlchemy} könyvtáron keresztül valósítottam meg. Az SQLAlchemy egy ORM\footnote{ORM: Object-Relational Mapping, azaz objektum-relációs leképezés. Az adatbázis tábláit a programkódban osztályokként, a sorokat pedig ezek példányaiként lehet kezelni.} réteget biztosít: az adatbázis táblái Python osztályokként (modellekként) definiálhatók, a lekérdezések pedig Python kifejezésekként fogalmazhatók meg. A típusannotált modelleknek köszönhetően a fejlesztőkörnyezet és a statikus típusellenőrzők már a kód írásakor jelzik a hibás mezőneveket és a típuseltéréseket. A könyvtár emellett paraméterezett lekérdezéseket állít elő, vagyis a paraméterek nem szövegösszefűzéssel kerülnek a lekérdezésbe, ami védelmet nyújt az SQL injection típusú támadásokkal szemben.

Az SQLAlchemy 2.0-s verziója natívan támogatja az aszinkron működést, ami a FastAPI aszinkron endpointjaival együtt lehetővé teszi, hogy a szerver egy adatbázis-művelet várakozási ideje alatt más kéréseket szolgáljon ki. Azokban az esetekben, ahol az ORM által generált lekérdezés nem lenne optimális, az SQLAlchemy Core alacsonyabb szintű, de továbbra is paraméterezett lekérdezésépítést is  lehetővé tesz, nyers SQL használata nélkül.

\subsubsection{Alembic}

A projekt fejlesztése során az adatbázis séma folyamatosan változott: új táblák és oszlopok kerültek bele, régiek módosultak vagy törlődtek. Ezeket a változásokat az Alembic \cite{Alembic} segítségével kezeltem, ami az SQLAlchemy-hez készült migrációs eszköz. Az Alembic a modellosztályok és az aktuális adatbázis összehasonlításával automatikusan legenerálja a migrációs szkriptet (programot), amit utána még kézzel szükség esetén lehet javítani. Minden migráció egy külön Python fájl, amiben szerepel a változtatás végrehajtása (\texttt{upgrade}) és visszavonása (\texttt{downgrade}) is.

Ennek a megoldásnak az a nagy előnye, hogy a séma változásai is verziókövetés alá kerülnek a kóddal együtt, és nem kell kézzel SQL szkripteket futtatni minden egyes változtatás után.

\subsubsection{Pydantic}

A Pydantic \cite{Pydantic} egy típusannotációkon alapuló adatvalidációs könyvtár, amelyre a FastAPI kérés- és válaszkezelése is épül. Egy Pydantic modell (séma) deklaratívan írja le az adatstruktúra mezőit, azok típusát és a rájuk vonatkozó megszorításokat (pl.\ a termék ára nem lehet negatív vagy az email-címnek érvényes formátumúnak kell lennie). Ha egy beérkező kérés törzse nem felel meg a sémának, a FastAPI automatikusan 422-es (Unprocessable Entity) státuszkóddal válaszol, a válaszban pedig megjelöli a hibás mezőket. Ennek eredményeként az endpointok üzleti logikája már validált, a sémának megfelelő adatokkal dolgozik.

\subsection{Object Storage}

A felhasználók és a körök által feltöltött fájlok, elsősorban a termékekhez tartozó ikonok tárolása nem az adatbázisban, hanem egy úgynevezett object storage-ban (objektumtárolóban) valósult meg. Az object storage olyan tárolási modell, amelyben a fájlok (objektumok) egyedi kulccsal azonosítva, egy lapos szerkezetű tárolóban (bucket) helyezkednek el, hierarchikus könyvtárszerkezet nélkül.

A legelterjedtebb ilyen szolgáltatás az Amazon S3 \cite{S3}, amelynek API-ja mára mondhatni szabvánnyá vált. A rendszer egy S3-kompatibilis tárolót használ, így a szerver kódja nem kötődik egy konkrét szolgáltatóhoz, és a tárolási háttér szükség esetén kis ráfordítással lecserélhető.

A képek adatbázison kívüli tárolását az indokolta, hogy a nagyobb bináris adatok jelentősen növelnék az adatbázis méretét, lassítanák a mentéseket, és minden letöltés a szerver erőforrásait terhelné. Object storage használata esetén a kliens előre beállított URL-en keresztül közvetlenül a tárolóból töltheti le a fájlokat, így a képek kiszolgálása nem terheli a szervert.

\subsection{Representational State Transfer (REST)}

A kliens és a szerver közötti kommunikáció a REST \cite{REST} elveit követi. A REST nem konkrét protokoll vagy könyvtár, hanem architekturális stílus, amelyet Roy Fielding fogalmazott meg 2000-ben megjelent doktori disszertációjában \cite{Fielding}. Alapgondolata, hogy a szerver által kezelt adatok erőforrásokként jelennek meg, amelyeket egy-egy URL azonosít (pl.\ \texttt{/groups/\{id\}/items}), az erőforrásokon végzett műveleteket pedig a HTTP metódusok fejezik ki:

\begin{itemize}
	\item \texttt{GET} -- erőforrás lekérdezése,
	\item \texttt{POST} -- új erőforrás létrehozása,
	\item \texttt{PUT} / \texttt{PATCH} -- meglévő erőforrás teljes vagy részleges módosítása,
	\item \texttt{DELETE} -- erőforrás törlése.
\end{itemize}

További alapelv az állapotmentesség: a szerver nem tárol a klienshez tartozó munkamenet-állapotot, ezért minden kérésnek tartalmaznia kell a kiszolgálásához szükséges összes információt (például a felhasználót azonosító tokent). A művelet eredményét a szerver HTTP státuszkóddal jelzi (pl.\ 200 -- sikeres, 201 -- létrehozva, 404 -- nem található, 403 -- nincs jogosultság). Az adatcsere a kérések és válaszok törzsében JSON formátumban történik.

A REST-alapú megközelítés mellett szólt egyszerűsége és széles körű elterjedtsége, valamint az, hogy a FastAPI is erre a modellre épül, és ugyanazt az API-t a mobil- és a webes kliens egyaránt használhatja.

\subsection{JSON Web Token (JWT)}

A bejelentkezett felhasználók azonosítása JSON Web Tokenek (JWT) \cite{JWT} segítségével történik. A JWT az RFC 7519-ben definiált szabványos formátum, amely három, ponttal elválasztott részből áll:

\begin{itemize}
	\item \textbf{Header:} a token típusát és az aláíráshoz használt algoritmust tartalmazza,
	\item \textbf{Payload:} az úgynevezett claim-eket tartalmazza, például a felhasználó azonosítóját (\texttt{sub}) és a lejárati időt (\texttt{exp}),
	\item \textbf{Signature:} az előző két rész digitális aláírása, amely csak a szerver titkos kulcsával állítható elő.
\end{itemize}

A három rész Base64URL kódolású, tehát a tartalmuk nem titkosított, bárki által dekódolható, ezért a payload érzékeny adatot nem tartalmazhat. Az aláírás ugyanakkor garantálja a token integritását, mivel a tartalom bármilyen módosítása az aláírás ellenőrzésekor kiderül. Sikeres bejelentkezéskor a szerver tokent állít ki a felhasználónak, amelyet a kliens minden további HTTP kérés \texttt{Authorization} fejlécében (Header) elküld. A szervernek így nem szükséges kérésenként az adatbázisból lekérdeznie a munkamenetet, elegendő az aláírást és a lejárati időt ellenőriznie, ami összhangban van a REST állapotmentességi elvével.

A JWT egyik hátránya, hogy egy kiállított token a lejárata előtt nem vonható vissza egyszerűen. A rendszer ezt kétféle token alkalmazásával kezeli: a hozzáférési token (access token) rövid élettartamú, mellette a kliens egy hosszabb lejáratú frissítő tokent (refresh token) is kap, amellyel új hozzáférési token igényelhető.

\subsection{OAuth2.0}

\subsection{Docker}

\section{Kliens (Frontend)}

\subsection{React Native}

\subsection{Expo}

